\begin{thebibliography}{99}

% --- PPMI / DaT-SPECT cohort and biomarker references ---
\bibitem{brumm2023milestone} Brumm, M. C., Siderowf, A., Simuni, T. \textit{et al.} PPMI: a milestone-based strategy to monitor Parkinson's disease progression. \textit{J. Parkinsons Dis.} \textbf{13}, 899--916 (2023).

\bibitem{marek2018ppmi} Marek, K., Chowdhury, S., Siderowf, A. \textit{et al.} The Parkinson's Progression Markers Initiative (PPMI)---establishing a PD biomarker cohort. \textit{Ann. Clin. Transl. Neurol.} \textbf{5}, 1460--1477 (2018).

\bibitem{dzialas2025dat} Dzialas, V., Bischof, G. N., M\"ollenhoff, K., Drzezga, A. \& van Eimeren, T. Dopamine transporter imaging as an objective monitoring biomarker in Parkinson's disease. \textit{Ann. Neurol.} \textbf{98}, 120--135 (2025).

\bibitem{nandhagopal2009} Nandhagopal, R., Kuramoto, L., Schulzer, M. \textit{et al.} Longitudinal progression of sporadic Parkinson's disease: a multi-tracer positron emission tomography study. \textit{Brain} \textbf{132}, 2970--2979 (2009).

\bibitem{pagano2018datspect} Pagano, G., Niccolini, F. \& Politis, M. Imaging in Parkinson's disease. \textit{Clin. Med.} \textbf{16}, 371--375 (2016).

% --- Mechanistic / physics-informed / SciML foundational references ---
\bibitem{fearnleyLees1991} Fearnley, J. M. \& Lees, A. J. Ageing and Parkinson's disease: substantia nigra regional selectivity. \textit{Brain} \textbf{114}, 2283--2301 (1991).

\bibitem{rackauckas2020universal} Rackauckas, C., Ma, Y., Martensen, J. \textit{et al.} Universal differential equations for scientific machine learning. Preprint at \textit{arXiv} 2001.04385 (2020).

\bibitem{chen2018neuralode} Chen, R. T. Q., Rubanova, Y., Bettencourt, J. \& Duvenaud, D. K. Neural ordinary differential equations. \textit{Adv. Neural Inf. Process. Syst.} \textbf{31}, 6571--6583 (2018).

\bibitem{deRooij2025physiology} de~Rooij, M., van~der~Graaf, P. H. \textit{et al.} Universal differential equations for physiological systems: regularisation strategies for interpretable hybrid models. Preprint (2025). Reference implementation: \url{https://github.com/Computational-Biology-TUe/ude-regularization}.

\bibitem{karniadakis2021pinn} Karniadakis, G. E., Kevrekidis, I. G., Lu, L., Perdikaris, P., Wang, S. \& Yang, L. Physics-informed machine learning. \textit{Nat. Rev. Phys.} \textbf{3}, 422--440 (2021).

\bibitem{willard2022pinnsurvey} Willard, J., Jia, X., Xu, S., Steinbach, M. \& Kumar, V. Integrating scientific knowledge with machine learning for engineering and environmental systems. \textit{ACM Comput. Surv.} \textbf{55}, 1--37 (2022).

% --- PD progression modelling references ---
\bibitem{oxtoby2021ebm} Oxtoby, N. P., Leyland, L.-A., Aksman, L. M. \textit{et al.} Sequence of clinical and neurodegeneration events in Parkinson's disease progression. \textit{Brain} \textbf{144}, 975--988 (2021).

\bibitem{severson2021progression} Severson, K. A., Chahine, L. M., Smolensky, L. A. \textit{et al.} Discovery of Parkinson's disease states and disease progression modelling: a longitudinal data study using machine learning. \textit{Lancet Digit. Health} \textbf{3}, e555--e564 (2021).

\bibitem{raj2015network} Raj, A., LoCastro, E., Kuceyeski, A., Tosun, D., Relkin, N. \& Weiner, M. Network diffusion model of progression predicts longitudinal patterns of atrophy and metabolism in Alzheimer's disease. \textit{Cell Rep.} \textbf{10}, 359--369 (2015).

\bibitem{simuni2018longitudinal} Simuni, T., Long, J. D., Caspell-Garcia, C. \textit{et al.} Predictors of time to initiation of symptomatic therapy in early Parkinson's disease. \textit{Ann. Clin. Transl. Neurol.} \textbf{5}, 1480--1494 (2018).

\bibitem{simuni2024nsdiss} Simuni, T., Chahine, L. M., Poston, K. \textit{et al.} A biological definition of neuronal $\alpha$-synuclein disease: towards an integrated staging system for research. \textit{Lancet Neurol.} \textbf{23}, 178--190 (2024).

\bibitem{iwaki2019gba} Iwaki, H., Blauwendraat, C., Leonard, H. L. \textit{et al.} Genetic risk of Parkinson disease and progression: an analysis of 13 longitudinal cohorts. \textit{Neurol. Genet.} \textbf{5}, e348 (2019).

\bibitem{paszke2019pytorch} Paszke, A., Gross, S., Massa, F. \textit{et al.} PyTorch: an imperative style, high-performance deep learning library. \textit{Adv. Neural Inf. Process. Syst.} \textbf{32}, 8024--8035 (2019).

\bibitem{chen2021torchdiffeq} Chen, R. T. Q. torchdiffeq: differentiable ODE solvers with full GPU support and adjoint backpropagation. \textit{GitHub} (2021). Available at \url{https://github.com/rtqichen/torchdiffeq}.

\bibitem{goetz2008mdsupdrs} Goetz, C. G., Tilley, B. C., Shaftman, S. R. \textit{et al.} Movement Disorder Society-sponsored revision of the Unified Parkinson's Disease Rating Scale (MDS-UPDRS): scale presentation and clinimetric testing results. \textit{Mov. Disord.} \textbf{23}, 2129--2170 (2008).

% --- Optimisation / architecture references ---
\bibitem{kingma2014adam} Kingma, D. P. \& Ba, J. Adam: a method for stochastic optimisation. \textit{Int. Conf. Learn. Represent.} (2015).

\bibitem{cho2014gru} Cho, K., van Merri\"enboer, B., G\"ul\c{c}ehre, \c{C}. \textit{et al.} Learning phrase representations using RNN encoder--decoder for statistical machine translation. \textit{Proc. Empir. Methods Nat. Lang. Process.} 1724--1734 (2014).

% --- Self-citations (limited to 3, 14% of 21 total) ---
\bibitem{dupre2026paper3} Dupre, B. D. Graph-informed digital twins for predicting NSD-ISS stage transitions in Parkinson's disease. Submitted to \textit{IEEE J. Biomed. Health Inform.} (2026).

\bibitem{dupre2026paper10} Dupre, B. D. A bidirectional-ready mechanistic twin for Parkinson's disease: external validation, calibration-aware observation likelihood, and NASEM 2024 self-audit. Submitted to \textit{npj Parkinsons Dis.} (2026).

\bibitem{dupre2026dissertation} Dupre, B. D. \textit{Graph-informed multimodal methods for Parkinson's disease biological staging, progression, and digital-twin integration}. PhD dissertation, University of North Dakota (2026, in preparation).

\end{thebibliography}
